% ============================================================
\chapter{Couples de Variables Al\'eatoires}
\label{ch:couples}
% ============================================================

Jusqu'ici, nous avons \'etudi\'e des variables al\'eatoires isol\'ees. Mais dans le monde r\'eel, les grandeurs sont rarement ind\'ependantes~: la taille et le poids d'un individu sont corr\'el\'es, le rendement d'une action d\'epend du march\'e global, la temp\'erature et la pression atmosph\'erique \'evoluent conjointement. Pour comprendre ces d\'ependances, il faut passer de la loi d'une variable \`a la \emph{loi conjointe} de plusieurs variables. Francis Galton, dans les ann\'ees 1880, a \'et\'e le premier \`a \'etudier syst\'ematiquement la corr\'elation --- en mesurant la taille des parents et de leurs enfants. Karl Pearson a ensuite formalis\'e le coefficient de corr\'elation qui porte son nom. Ce chapitre d\'eveloppe les outils fondamentaux~: lois conjointes, marginales, conditionnelles, covariance et corr\'elation.

\begin{intuition}
Ce chapitre \'etudie les lois conjointes de deux (ou plusieurs) variables al\'eatoires,
les lois marginales et conditionnelles, la covariance et la corr\'elation. Ces outils
sont essentiels pour mod\'eliser les d\'ependances entre grandeurs al\'eatoires.
\end{intuition}

% ------------------------------------------------------------
\section{Loi conjointe}
% ------------------------------------------------------------

\subsection{Cas discret}

\begin{definition}[Loi conjointe discr\`ete]
Soit $(X, Y)$ un couple de variables al\'eatoires discr\`etes. La \textbf{loi
conjointe} est la fonction~:
\[
    p_{X,Y}(x_i, y_j) = \PP(X = x_i, Y = y_j), \quad
    \text{avec } \sum_{i,j} p_{X,Y}(x_i, y_j) = 1.
\]
\end{definition}

\subsection{Cas continu}

\begin{definition}[Densit\'e conjointe]
Le couple $(X, Y)$ admet une \textbf{densit\'e conjointe} $f_{X,Y}$ si pour tout
bor\'elien $B \subseteq \R^2$~:
\[
    \PP((X,Y) \in B) = \iint_B f_{X,Y}(x, y)\, dx\, dy,
\]
avec $f_{X,Y}(x,y) \geq 0$ et $\iint_{\R^2} f_{X,Y}(x,y)\, dx\, dy = 1$.
\end{definition}

\begin{definition}[Fonction de r\'epartition conjointe]
\[
    F_{X,Y}(x, y) = \PP(X \leq x, Y \leq y).
\]
Dans le cas \`a densit\'e~: $f_{X,Y}(x,y) = \frac{\partial^2 F_{X,Y}}{\partial x\, \partial y}(x,y)$.
\end{definition}

% ------------------------------------------------------------
\section{Lois marginales}
% ------------------------------------------------------------

\begin{definition}[Marginales]
Les \textbf{lois marginales} s'obtiennent par sommation (discret) ou int\'egration
(continu) sur l'autre variable~:
\begin{align*}
    &\text{Discret~:} & p_X(x_i) &= \sum_j p_{X,Y}(x_i, y_j), &
    p_Y(y_j) &= \sum_i p_{X,Y}(x_i, y_j). \\
    &\text{Continu~:} & f_X(x) &= \int_{-\infty}^{+\infty} f_{X,Y}(x, y)\, dy, &
    f_Y(y) &= \int_{-\infty}^{+\infty} f_{X,Y}(x, y)\, dx.
\end{align*}
\end{definition}

\begin{attention}[title=\textbf{La loi conjointe d\'etermine les marginales, mais pas l'inverse}]
Connaissant $f_{X,Y}$, on calcule $f_X$ et $f_Y$. Mais connaissant seulement
$f_X$ et $f_Y$, on ne peut pas reconstituer $f_{X,Y}$ sans information sur la
d\'ependance entre $X$ et $Y$.
\end{attention}

\begin{exemple}
Soit $(X,Y)$ de densit\'e conjointe $f_{X,Y}(x,y) = 6(1-y)$ pour
$0 < x < y < 1$ et $0$ sinon. Alors~:
\[
    f_X(x) = \int_x^1 6(1-y)\, dy = 3(1-x)^2, \quad 0 < x < 1.
\]
\[
    f_Y(y) = \int_0^y 6(1-y)\, dx = 6y(1-y), \quad 0 < y < 1.
\]
\end{exemple}

% ------------------------------------------------------------
\section{Loi conditionnelle}
% ------------------------------------------------------------

\begin{definition}[Densit\'e conditionnelle]
La \textbf{densit\'e conditionnelle} de $Y$ sachant $X = x$ est~:
\[
    f_{Y \mid X}(y \mid x) = \frac{f_{X,Y}(x, y)}{f_X(x)}, \quad f_X(x) > 0.
\]
De m\^eme dans le cas discret~: $\PP(Y = y_j \mid X = x_i) = p_{X,Y}(x_i, y_j)/p_X(x_i)$.
\end{definition}

% ------------------------------------------------------------
\section{Ind\'ependance}
% ------------------------------------------------------------

\begin{theoreme}[Caract\'erisation de l'ind\'ependance]
$X$ et $Y$ sont ind\'ependantes si et seulement si~:
\begin{itemize}
    \item \textbf{Cas discret~:} $p_{X,Y}(x,y) = p_X(x)\, p_Y(y)$ pour tout $(x,y)$.
    \item \textbf{Cas continu~:} $f_{X,Y}(x,y) = f_X(x)\, f_Y(y)$ pour tout $(x,y)$.
    \item \textbf{Via la CDF~:} $F_{X,Y}(x,y) = F_X(x)\, F_Y(y)$ pour tout $(x,y)$.
\end{itemize}
\end{theoreme}

\begin{proposition}[Crit\`ere de factorisation]
Le couple $(X,Y)$ continu est ind\'ependant si et seulement si il existe des
fonctions $g, h$ telles que $f_{X,Y}(x,y) = g(x)\, h(y)$ sur le support (qui
doit \^etre un produit cart\'esien).
\end{proposition}

% ------------------------------------------------------------
\section{Covariance et corr\'elation}
% ------------------------------------------------------------

\begin{definition}[Covariance]
\[
    \Cov(X, Y) = \E[(X - \E[X])(Y - \E[Y])] = \E[XY] - \E[X]\E[Y].
\]
\end{definition}

\begin{proposition}[Propri\'et\'es de la covariance]
\begin{enumerate}[label=(\roman*)]
    \item $\Cov(X, X) = \Var(X)$.
    \item $\Cov(X, Y) = \Cov(Y, X)$ (sym\'etrie).
    \item $\Cov(aX + b, cY + d) = ac\, \Cov(X, Y)$ (bilin\'earit\'e).
    \item $\Var(X + Y) = \Var(X) + \Var(Y) + 2\Cov(X, Y)$.
    \item Si $X \indep Y$, alors $\Cov(X, Y) = 0$.
    \item $\Var\!\left(\sum_{i=1}^n X_i\right) = \sum_{i=1}^n \Var(X_i)
          + 2\sum_{i < j} \Cov(X_i, X_j)$.
\end{enumerate}
\end{proposition}

\begin{definition}[Coefficient de corr\'elation]
\[
    \rho(X, Y) = \frac{\Cov(X, Y)}{\sqrt{\Var(X)\, \Var(Y)}}.
\]
\end{definition}

\begin{theoreme}[Propri\'et\'es de $\rho$]
\begin{enumerate}[label=(\roman*)]
    \item $-1 \leq \rho(X, Y) \leq 1$ \quad (cons\'equence de Cauchy--Schwarz).
    \item $\abs{\rho(X,Y)} = 1$ si et seulement si $Y = aX + b$ p.s.\ avec
          $a \neq 0$.
    \item $\rho(X,Y) = 0$ signifie que $X$ et $Y$ sont \textbf{non corr\'el\'ees}.
\end{enumerate}
\end{theoreme}

\begin{attention}[title=\textbf{Non corr\'el\'e $\neq$ ind\'ependant}]
$\Cov(X,Y) = 0$ n'implique pas l'ind\'ependance. Exemple~: si $X \sim \mathcal{N}(0,1)$
et $Y = X^2$, alors $\Cov(X, Y) = \E[X^3] = 0$ (par sym\'etrie), mais $X$ et $Y$
ne sont \'evidemment pas ind\'ependants.
\end{attention}

% ------------------------------------------------------------
\section{Loi normale bivari\'ee}
% ------------------------------------------------------------

\begin{definition}[Normale bivari\'ee]
Le vecteur $(X, Y)$ suit une \textbf{loi normale bivari\'ee}
$\mathcal{N}_2(\boldsymbol{\mu}, \Sigma)$ avec
$\boldsymbol{\mu} = (\mu_X, \mu_Y)^T$ et
$\Sigma = \begin{pmatrix} \sigma_X^2 & \rho\sigma_X\sigma_Y \\
\rho\sigma_X\sigma_Y & \sigma_Y^2 \end{pmatrix}$ si~:
\[
    f_{X,Y}(x,y) = \frac{1}{2\pi\sigma_X\sigma_Y\sqrt{1-\rho^2}}
    \exp\!\left(-\frac{1}{2(1-\rho^2)}\left[
    \frac{(x-\mu_X)^2}{\sigma_X^2}
    - 2\rho\frac{(x-\mu_X)(y-\mu_Y)}{\sigma_X\sigma_Y}
    + \frac{(y-\mu_Y)^2}{\sigma_Y^2}
    \right]\right).
\]
\end{definition}

\begin{theoreme}[Propri\'et\'es de la normale bivari\'ee]
\begin{enumerate}[label=(\roman*)]
    \item Les marginales sont $X \sim \mathcal{N}(\mu_X, \sigma_X^2)$ et
          $Y \sim \mathcal{N}(\mu_Y, \sigma_Y^2)$.
    \item $\rho(X,Y) = \rho$ (le param\`etre de la loi).
    \item $X$ et $Y$ sont ind\'ependantes si et seulement si $\rho = 0$.
    \item Toute combinaison lin\'eaire $aX + bY$ est normale.
\end{enumerate}
\end{theoreme}

\begin{remarque}
Le point (iii) est sp\'ecifique \`a la loi normale~: en g\'en\'eral, non corr\'el\'e
n'implique pas ind\'ependant, mais pour un vecteur gaussien, si.
\end{remarque}

% ------------------------------------------------------------
\section{Somme de variables al\'eatoires}
% ------------------------------------------------------------

\begin{theoreme}[Convolution]
Si $X$ et $Y$ sont ind\'ependantes avec densit\'es $f_X$ et $f_Y$, alors
$Z = X + Y$ a pour densit\'e la \textbf{convolution}~:
\[
    f_Z(z) = (f_X * f_Y)(z) = \int_{-\infty}^{+\infty} f_X(t)\, f_Y(z - t)\, dt.
\]
\end{theoreme}

\begin{exemple}
Si $X \sim \text{Exp}(\lambda)$ et $Y \sim \text{Exp}(\lambda)$ sont ind\'ependantes,
alors pour $z > 0$~:
\[
    f_Z(z) = \int_0^z \lambda e^{-\lambda t}\, \lambda e^{-\lambda(z-t)}\, dt
    = \lambda^2 e^{-\lambda z} \int_0^z dt = \lambda^2 z\, e^{-\lambda z}.
\]
C'est la densit\'e de $\Gamma(2, \lambda)$.
\end{exemple}

% ------------------------------------------------------------
\section{Exercices}
% ------------------------------------------------------------

\begin{exercice}
Soit $(X,Y)$ de densit\'e $f(x,y) = c\, x\, y$ pour $0 < x < 1$, $0 < y < 1$
et 0 sinon. D\'eterminer $c$, les marginales, et v\'erifier l'ind\'ependance.
\end{exercice}

\begin{exercice}
Soit $(X,Y)$ de densit\'e $f(x,y) = 2$ pour $0 < x < y < 1$, $0$ sinon.
Calculer $f_X$, $f_Y$, $\E[X]$, $\E[Y]$, $\Cov(X,Y)$, $\rho(X,Y)$.
\end{exercice}

\begin{exercice}
Montrer que si $(X,Y) \sim \mathcal{N}_2$ avec $\rho = 0$, alors $X$ et $Y$
sont ind\'ependantes (factoriser la densit\'e conjointe).
\end{exercice}

\begin{exercice}
Soit $X \sim \mathcal{N}(0,1)$ et $Y = X^2$. V\'erifier que $\Cov(X,Y) = 0$
mais que $X$ et $Y$ ne sont pas ind\'ependantes.
\end{exercice}

\begin{exercice}
Calculer par convolution la densit\'e de $Z = X + Y$ o\`u
$X \sim \mathcal{N}(0,1)$ et $Y \sim \mathcal{N}(0,1)$ sont ind\'ependantes.
\end{exercice}

\begin{exercice}
Soit $(X,Y)$ uniforme sur le disque $\{(x,y) : x^2 + y^2 \leq 1\}$.
Calculer les marginales, $\E[X]$, $\Var(X)$, $\Cov(X,Y)$.
\end{exercice}

\begin{exercice}
On tire au hasard un point $(X,Y)$ uniform\'ement dans le triangle de sommets
$(0,0)$, $(1,0)$, $(0,1)$. D\'eterminer la densit\'e conjointe, les marginales,
et calculer $\rho(X,Y)$.
\end{exercice}

\begin{exercice}
Soit $X_1, X_2, X_3$ i.i.d.\ $\text{Exp}(1)$. Calculer $\Cov(X_1 + X_2, X_2 + X_3)$.
\end{exercice}
